% ====================================================================
% ch1v22_partT_divider.tex — Ch1 Part T(흑연 열특성) 파트 구분면 (R1 신설)
%   R2 확정(MISC §C): 다리 문단 확정 + 폭 서식 live 참조.
% ====================================================================
\clearpage
\begin{center}
{\LARGE\bfseries Part T --- 흑연 열특성: 엔트로피 분해와 가역 발열}
\end{center}
\addcontentsline{toc}{section}{Part T --- 흑연 열특성: 엔트로피 분해와 가역 발열}
\vspace{0.6em}
여기까지가 흑연 $\dd Q/\dd V$ 한 곡선을 그리는 계산 사슬이었다. 이제 같은 전극$\cdot$같은 입력에서
곡선 대신 \emph{열}을 뽑는다 --- 전이별 반응 엔트로피 $\Delta S_{\rxn,j}$ 가 어디서 오는지(배치$\cdot$진동$\cdot$전자
세 분포), 그리고 그것이 가역 발열 $\dot Q_\rev=-IT\,\partial U_\oc/\partial T$ 로 어떻게 조립되는지다.
앞 파트의 전하 보존 반전~\eqref{eq:sm-mc-balance}$\cdot$평형 종~\eqref{eq:eqpeak}$\cdot$폭 서식(\S\ref{sec:width})이 그대로 입력으로 들어온다.
